% xelatex script.tex
\documentclass[12pt,a4paper]{article}

\usepackage{fontspec}
\setmainfont{Times New Roman}
\usepackage[margin=2.2cm]{geometry}
\usepackage{xcolor}
\usepackage{enumitem}
\setlength{\parindent}{0pt}
\setlength{\parskip}{0pt}
\pagestyle{empty}

% перевод под репликами: \translationfalse, чтобы убрать
\newif\iftranslation
\translationtrue

\newcommand{\rem}[1]{\par\medskip{\itshape\color{black!65}#1}\par\medskip}
\newcommand{\tr}[1]{\iftranslation\\{\small\itshape\color{black!55}#1}\fi}
\newcommand{\say}[3]{\par\smallskip\textbf{#1:}\ #2\tr{#3}\par}

\newcommand{\RU}{РУСЛАН}
\newcommand{\KR}{КРИСТИАН}
\newcommand{\SO}{СОЛОМИЯ}
\newcommand{\FI}{ФИДАН}

\begin{document}

\begin{center}
{\LARGE\bfseries Eksam — null}\\[4pt]
\end{center}

\section*{Логлайн}

После проваленного экзамена по эстонскому четверо незнакомцев по одному подсаживаются на одну лавочку. Они разбирают, где завалились, и заедают провал тем, что напоминает о доме. Единственное, от чего все дружно отказываются, — «чай» тихой Фидан.

\section*{Персонажи}

\begin{itemize}[leftmargin=*]
\item \textbf{Руслан} (Казахстан) — валит экзамен не в первый раз и давно привык. Не интроверт, просто его уже ничем не удивишь: спокоен, общается легко. Утешает новичков как бывалый, и от этих утешений им только хуже.
\item \textbf{Кристиан} (Латвия) — драма-квин, всё латвийское лучше. Страдает громче всех, хотя до дома 4 часа на автобусе.
\item \textbf{Соломия} (Украина) — отличница, не добрала один балл. Не может это отпустить и всем доказывает, что экзамен был несправедливым.
\item \textbf{Фидан} (Азербайджан) — интроверт из большой заботливой семьи. Телефон звонит не переставая, и каждый раз она тихо отвечает «Hə, ana. Çay içirəm» («Да, мам. Чай пью») и отпивает из термоса. Из вежливости предлагает всем, все отказываются. В финале кто-то наконец соглашается, отпивает и замирает: в термосе коньяк.
\end{itemize}

\section*{Стиль}

Ситком. Снимаем одним планом: камера на штативе, общий план на лавочку, все четверо всё время в кадре. Играем как в театре: реплика, пауза, реакция.

\rem{Первый кадр: пустая лавочка, титр с названием «Eksam — null».}

\section*{Ругательства}

Каждый, когда эмоции зашкаливают, срывается на родной язык. Ругаемся без мата, остальные не понимают и вежливо кивают.

\begin{itemize}[leftmargin=*,nosep]
\item \textbf{Руслан}: «Ойбай», «Әй, құдай-ау», «Құдай сақтасын»
\item \textbf{Кристиан}: «Sasodīts», «Pie velna», «Velns»
\item \textbf{Соломия}: «Трясця його матері», «Холера», «Щоб вас усіх»
\item \textbf{Фидан}: «Vay dədə vay», «Ay Allah». Она ругается всего один раз, и от неё это звучит неожиданно.
\end{itemize}

\section*{Действия}

\subsection*{Акт 1. Подсаживаются по одному}
\begin{enumerate}[leftmargin=*]
\item Пустая лавочка, титр.
\item Фидан садится одна с термосом. Звонит телефон: «Hə, ana. Çay içirəm». Отпивает.
\item Руслан спокойно садится рядом, у него кефир. Видит её лист с результатами: она тоже завалила. Рассказывает, что сдаёт не первый раз и ему норм. Фидан предлагает термос, Руслан отказывается: у него кефир.
\item Кристиан подходит с драмой и ругается по-латышски. Садится и страдает, как будто жизнь кончена. Руслан его утешает, и становится только хуже. Фидан предлагает термос, Кристиан отказывается.
\item Соломия влетает в ярости и ругается по-украински: 59 баллов из 60. Фидан предлагает термос, Соломия отказывается.
\item Сравнивают, кто на чём завалился. У Фидан ноль за говорение: на устной части она не сказала ни слова.
\end{enumerate}

\subsection*{Акт 2. Еда и дом}
\begin{enumerate}[leftmargin=*]
\item Все достают свою еду. У Соломии kohuke: она заедает стресс, он похож на сырники. У Кристиана эстонский хлеб, хвалит его и тут же поправляется, что латвийский лучше. У Руслана кефир.
\item Пробуют еду друг друга. У каждого своя реакция, кто-то срывается на родной язык.
\item Фидан звонят второй раз: «Çay içirəm», глоток. Снова предлагает термос, все снова отказываются.
\item Говорят, по чему скучают. Неожиданно ломается спокойный Руслан: заговорил о маминой еде, раскис, «Ойбай…». Остальные утешают его своей едой.
\end{enumerate}

\subsection*{Акт 3. Тост}
\begin{enumerate}[leftmargin=*]
\item Следующий экзамен в декабре, настроение на нуле. Кто-то предлагает выпить, «Terviseks», но выпить нечего: кефир кончился, чая ни у кого нет.
\item Остаётся только термос Фидан. Третий звонок. Перед тем как ответить, Фидан впервые ругается: «Vay dədə vay…». Потом: «Hə, ana. Çay içirəm». Глоток долгий.
\item Руслан сдаётся: «Okei. Tee». Отпивает, замирает и смотрит на Фидан. Она спокойно пожимает плечами.
\item Остальные молча тянут к ней стаканчики.
\item Финал: у Фидан снова звонит телефон, и все четверо хором отвечают: «Hə, ana. Çay içirəm».
\end{enumerate}

\end{document}
